% Body fragment: the including document supplies amsmath and amssymb.
\section{Rigid blocks and comparative absolute error}
\label{sec:rigid-theory}

The rigid-component construction is classical. Theorems~6--7 of
Tsamardinos, Muscettola, and Morris~\cite{tsamardinos1998} directly precede
the SCC characterization and representative/offset substitution below.
Johnson reweighting is likewise classical~\cite{johnson1977}.
We give a self-contained specialization to the declared loss and its
reduction parameter; this attribution is not a novelty certificate.

Let \(I=\{1,\ldots,n\}\), \(I_0=I\cup\{0\}\), and \(t_0=0\).
Assume the finite directed constraints \(t_j-t_i\le b_{ij}\) define a
nonempty bounded set \(S\). All inequalities are non-strict, so \(S\) is
compact. In the arc orientation \(i\to j\), let \(D_{ij}\) be the
shortest-path distance. These distances are finite and satisfy
\[
 D_{ij}=\max_{t\in S}(t_j-t_i),\qquad
 \ell_i=-D_{i0},\quad u_i=D_{0i}.
\]
Indeed every path bounds the difference, while
\(t_j=D_{pj}-D_{p0}\) is an anchored feasible vector attaining each
bound from \(p\). The triangle inequality makes both \(\ell\) and \(u\)
feasible. These are exact marginal endpoints, not independent box
constraints replacing the STN.

For finite fixed \(a_i,c_i\) and \(w_i>0\), put
\[
 f_i(z)=|c_i-z|-|a_i-z|,\quad
 F(t)=\sum_{i\in I}w_i f_i(t_i),\quad
 \Lambda=\sum_{i\in I}w_i,\quad
 W=\max_{t\in S}F(t).
\]
Coordinate 0 has no loss; the comparative MAE is \(W/\Lambda\).
Each unary term is globally monotone. Let \(p_{\rm orig},n_{\rm orig}\)
count its effectively nonconstant increasing and decreasing coordinates,
respectively, and \(r_{\rm orig}=\min(p_{\rm orig},n_{\rm orig})\).
Effective constancy is exactly \(f_i(\ell_i)=f_i(u_i)\).

\subsection{Classical rigid equivalence and offsets}

Choose any feasible potential \(h\), normalized by \(h_0=0\), and form
\(b'_{ij}=b_{ij}+h_i-h_j\ge0\). This arbitrary potential need not be the
forward-0 distance profile. Consider the directed subgraph of exactly
zero reduced-cost arcs.

\noindent\textbf{Lemma 1 (rigid equivalence).}
Its SCCs are precisely the maximal classes for which \(t_j-t_i\) is
constant over \(S\); the constant is \(h_j-h_i\).

\noindent\emph{Proof.}
Zero-cost paths in both directions give the two opposite bounds
forcing equality. Conversely a constant difference has that value
because \(h\in S\). Tightness gives
\(D_{ij}=h_j-h_i\) and \(D_{ji}=h_i-h_j\).
Both reduced shortest-path distances are zero. Nonnegative costs force
zero-length paths to consist entirely of zero-cost arcs.
Thus the vertices lie in one zero-cost SCC. The constant-difference
definition also proves independence from the chosen potential.
\(\square\)

Select a representative \(q_B\) in each block \(B\), with \(q_{B_0}=0\)
for the origin block. Set
\[
 \delta_i=h_i-h_{q_B},\qquad t_i=x_B+\delta_i
 \quad(i\in B).
\]
An arc \(i\in B\to j\in C\) becomes
\[
 x_C-x_B\le b_{ij}+\delta_i-\delta_j.                 \tag{1}
\]
These constraints with \(x_{B_0}=0\) define an exact quotient:
every original truth satisfies them by Lemma~1, and every quotient
truth reconstructs an original truth satisfying every arc by
substitution. Representatives make the correspondence one-to-one.
Internal arcs are automatic; parallel arcs may be replaced by their
minimum bound while retaining original identities.
Consequently
\[
 \begin{aligned}
 D^Q_{BC}&=D_{q_Bq_C},\quad L_B=\ell_{q_B},\quad U_B=u_{q_B},\\
 \ell_i&=L_B+\delta_i,\qquad u_i=U_B+\delta_i .
 \end{aligned}
\]
No argument requires the original truth set to have full dimension.
A single zero-cost arc, or weak connectivity of zero arcs, does not
establish rigid equivalence.

\subsection{Aggregate slopes and the reduction parameter}

For \(B\ne B_0\), define
\[
 A_i=a_i-\delta_i,\quad C_i=c_i-\delta_i,\qquad
 g_B(x)=\sum_{i\in B}w_i\bigl(|C_i-x|-|A_i-x|\bigr).
\]
All members of \(B_0\) are fixed. With
\[
 C_0=\sum_{i\in B_0\cap I}w_i
       \bigl(|c_i-\delta_i|-|a_i-\delta_i|\bigr),
 \qquad F(t)=C_0+\sum_{B\ne B_0}g_B(x_B),             \tag{2}
\]
arbitrary offsets and different unary pairs are preserved exactly.
They need not combine into one weighted absolute-difference term.

On each open segment between shifted kinks, the aggregate slope is
\[
 s_B(x)=\sum_{i\in B}2w_i\operatorname{sign}(a_i-c_i)
  \mathbf{1}_{\{\min(A_i,C_i)<x<\max(A_i,C_i)\}}.     \tag{3}
\]
Inspect all positive-length segments in the exact interval
\([L_B,U_B]\). Let \(\mathcal Z\) contain flat blocks, including
singleton marginals; let \(\mathcal P\) contain nonconstant blocks
whose segment slopes are all nonnegative, and \(\mathcal N\) those
whose slopes are all nonpositive. The remaining blocks
\(\mathcal B\) have slopes of both signs. These cases partition the
nonorigin blocks. Sorting and combining the at most \(2n\) original
events gives the classification in \(O(n\log(n+1))\) arithmetic and
comparison operations.

Endpoint comparison is insufficient: for \(t_1=t_2=x\in[0,2]\),
\(a=(1,1)\), \(c=(0,2)\), and unit weights,
\[
 g(x)=
 \begin{cases}2x,&0\le x\le1,\\4-2x,&1\le x\le2.\end{cases}
\]
Thus \(g(0)=g(2)=0\), whereas \(g(1)=2\).
Flat block values may be nonzero: retain them in (2), or absorb them
into \(C_0\). In either case \(\Lambda\) remains the total original
weight, including all fixed and flat coordinates.

\noindent\textbf{Lemma 2 (parameter domination).}
The aggregate parameter satisfies
\[
 r_\star=|\mathcal B|+\min(|\mathcal P|,|\mathcal N|),
 \qquad r_\star\le r_{\rm orig}.                    \tag{4}
\]

\noindent\emph{Proof.}
Every nonconstant \(\mathcal P\) block contains an effectively
increasing original term, every \(\mathcal N\) block a decreasing
term, and every \(\mathcal B\) block at least one of each.
Otherwise its sum could not have the declared slopes on its exact
marginal. Offsets preserve direction and effective activity.
Blocks are disjoint, so
\[
 |\mathcal B|+|\mathcal P|\le p_{\rm orig},\qquad
 |\mathcal B|+|\mathcal N|\le n_{\rm orig}.
\]
Taking the smaller bound proves (4). \(\square\)

\subsection{Exact conditional representative reduction}

For any selected nonorigin blocks \(K\), define
\[
 Q_K=\{z:L_C\le z_C\le U_C,\
              z_E-z_C\le D^Q_{CE}\ (C,E\in K)\}.
\]
Adjoin \(z_{B_0}=0\), write \(K_0=K\cup\{B_0\}\), and set
\[
 m_B(z)=\max_{C\in K_0}(z_C-D^Q_{BC}),\qquad
 M_B(z)=\min_{C\in K_0}(z_C+D^Q_{CB}).               \tag{5}
\]
For empty \(K\) these are \(L\) and \(U\).

\noindent\textbf{Lemma 3 (conditional extrema).}
\(Q_K\) is the exact projection of the quotient STN.
For every \(z\in Q_K\), both extensions in (5) are feasible, agree
with \(z\) on \(K\), and bound every matching feasible \(x\) as
\(m(z)\le x\le M(z)\).

\noindent\emph{Proof.}
Projected inequalities bound all branches at a selected block, and
its own branch gives equality. The triangle inequality yields
\[
 z_C-D^Q_{EC}\le z_C-D^Q_{BC}+D^Q_{BE}
                   \le m_B+D^Q_{BE};
\]
taking the maximum proves feasibility of \(m\).
If \(C\) attains \(M_B\), then
\(M_E\le z_C+D^Q_{CE}\le M_B+D^Q_{BE}\), proving the other case.
Any matching feasible \(x\) lies above each lower branch and below
each upper branch. Necessity of \(Q_K\) is immediate; the extensions
prove sufficiency. This also handles lower-dimensional domains.
\(\square\)

\noindent\textbf{Theorem 1 (block-minority reduction).}
If \(|\mathcal P|\le|\mathcal N|\), choose
\(K=\mathcal B\cup\mathcal P\). Then
\[
 W=\max_{z\in Q_K}\left[C_0+
                  \sum_{B\ne B_0}g_B(m_B(z))\right].              \tag{6}
\]
For \(|\mathcal N|\le|\mathcal P|\), choose
\(K=\mathcal B\cup\mathcal N\) and use \(M\).

\noindent\emph{Proof.}
In the first case, all nonmonotone and increasing block losses are
fixed on a fibre. Remaining nonconstant losses are nonincreasing,
so the least extension maximizes their sum on that fibre.
It is feasible, and the projection is exact; maximizing proves
(6). The greatest-extension case is symmetric. Compactness gives
attainment, and offsets reconstruct an original worst truth.
\(\square\)

For \(r_\star=0\) the appropriate endpoint extension suffices,
even when original individual directions are mixed.
This parameter is a sufficient reduction size, not a claim of
minimal possible dimension.

\subsection{One representative and strict separation families}

For \(r_\star=1\), let \(p\) be the selected representative.
Using original distances and original unary terms, the path is
\[
 t_i(z)=\max(\ell_i,z-D_{ip})
 \quad\hbox{or}\quad
 t_i(z)=\min(u_i,z+D_{pi}),\qquad z\in[\ell_p,u_p].   \tag{7}
\]
These formulas enforce all offsets automatically. They remain
valid when the selected block loss is nonmonotone.
Put \(l_i=\min(a_i,c_i)\), \(v_i=\max(a_i,c_i)\).
On the least path, a nonconstant term has slope
\(s_i=2w_i\operatorname{sign}(a_i-c_i)\) between
\[
 \max(\ell_i,l_i)+D_{ip}\quad\hbox{and}\quad v_i+D_{ip};
\]
skip \(a_i=c_i\) or \(\ell_i\ge v_i\).
On the greatest path the event interval is
\[
 \bigl(l_i-D_{pi},\,\min(u_i,v_i)-D_{pi}\bigr);
\]
skip \(a_i=c_i\) or \(u_i\le l_i\).
Each original term has at most two events. Between consecutive
events the objective is affine, so an endpoint or event attains
the maximum: at most \(2n+2\) candidates before deduplication.
Initialize the right slope at the left domain endpoint, combine
simultaneous events, and evaluate each event before updating its
slope. No grid is required. Given the required profiles, the
mathematical sweep takes \(O(n\log(n+1))\) operations and \(O(n)\)
extra storage. Checking its witness against all \(M\) original arcs
gives \(O(M+n\log(n+1))\) postprocessing, excluding preprocessing.

Two analytical families give strict separation from the original
coordinate parameter. Let \(n=2m\), all truths equal
\(x\in[0,2]\), and all weights be one.
Repeating \(m\) pairs \(a=(1,0)\), \(c=(0,1)\) cancels identically:
\[
 r_{\rm orig}=m,\qquad r_\star=0,\qquad W=W/\Lambda=0.
\]
Repeating \(m\) pairs \(a=(1,1)\), \(c=(0,2)\) gives the peak above:
\[
 r_{\rm orig}=m,\quad r_\star=1,\quad
 W=2m\ \hbox{at }x=1,\quad \Lambda=2m,\quad W/\Lambda=1.
\]
Arbitrary rigid offsets preserve these identities when the same
offset is added to each truth, baseline, and candidate entry.
The displayed values assume unit weights; arbitrary positive
weights are permitted by the theorem but need not preserve
cancellation. This compares with the unreduced coordinate
parameter, not every classical presolver, and reports no experiment.

Interior alternatives use pairs \(a=(3/2,1/2)\), \(c=(1/2,3/2)\)
for cancellation, or \(a=(1,1)\), \(c=(1/2,3/2)\) for the peak.
The latter pair is zero outside \([1/2,3/2]\), rises with slope 2
to height 1 at \(x=1\), then falls with slope \(-2\).
Thus the respective \(r_\star\) values remain 0 and 1; the peak
family has \(W=m\) and \(W/\Lambda=1/2\).
Both alternatives have \(k=\#\{i:\ell_i<c_i<u_i\}=n\).
This separates the parameters, not LP running times or lower bounds:
common-truth families have few distinct levels and feasible sign patterns.

\subsection{General dimension, implementation scope, and checking}

For \(r=r_\star\ge1\), partition \(Q_K\) by equalities between
conditional affine branches, their crossings of every shifted
original unary kink, and projected-domain boundaries.
Discard constant equalities and combine duplicate hyperplanes.
With \(b\) blocks, there are
\(H=O(br^2+nr+r^2)=O(nr^2)\) hyperplanes.
The objective is affine on each cell closure, so some optimum is a
cell vertex. Enumerating \(r\) independent active hyperplanes,
checking the projection, and evaluating the extension gives
\[
 O\!\left(H^r(r^3+r^2+nr)\right)
\]
arithmetic complexity. Domain equalities supply missing normals
in lower-dimensional cases. This is XP for fixed \(r\), not FPT
or general polynomial time. Incomplete enumeration does not certify
an optimum.

Distance profiles in the intended implementation remain on the
\emph{original graph}. The quotient supports classification,
representative selection, and the equivalence proof; no fully
contracted shortest-path solver is presumed. Classical feasible
potentials and forward/reverse Dijkstra require \(2r+2\) profiles.
Reverse the already reweighted graph, and retain original identities
when coalescing parallel arcs. With \(M\) original input arcs, a
safe preprocessing and classification bound is
\[
 O\!\left(nM+(r+1)(M+n)\log(n+1)+n\log(n+1)\right),
\]
with \(O(M+(r+1)n)\) storage. The Bellman--Ford term remains.

Independent checking is a separate obligation: validate input
binding, potential, SCCs, offsets, every aggregate slope segment
and constant, selected representatives, required distance profiles,
complete sweep or arrangement coverage, original feasibility,
objective, and original normalization. A feasible witness alone
gives only a lower bound on \(W\). No checker execution or pass is
asserted here. Exact arithmetic bounds do not remove rational bit
costs; real inputs require an exact-real interpretation.
Floating-point zero tests do not certify rigid contraction.
The result establishes no new shortest-path method, general
polynomial or FPT algorithm, physical accuracy, or publication-level
novelty.
